% Appendix (v3): a formal account of the three trade-offs of Section 6. Definitions follow
% Section 3; nothing here is a new measurement. Every number is a \res key or a count already
% printed in the tables it cites.
\section{A Formal Account of the Trade-off}
\label{app:formal}

The results of Section~\ref{sec:results} share one cause: a claim, not its passage, is the
generator's input. This appendix states the three consequences precisely. The statements are
elementary; their use is to say which observed differences are structural, so that they will
recur in any system built this way, and which depend on a particular build.

\paragraph{Notation.}
A stored claim is $c=(s,p,o,\mathrm{vol}(c),\mu(c))$: subject, predicate, object or value,
the volume of its first evidence (\firstvol) and a polarity
$\mu(c)\in\{\textit{asserted},\textit{presumed}\}$. The implemented schema stores no speaker
and no finer modality (prediction, hypothesis, belief). For an entity $e$, let
$I(e)$ be its verifier-accepted intervals and $I_t(e)=\{i\in I(e):\mathrm{start}(i)\le t\}$
those eligible at cutoff $t$; the volume filter is $\phi_t(c)=\mathbf{1}[\mathrm{vol}(c)\le t]$.
$\Gt(e)$ denotes the gold facts about $e$ eligible at $t$ (Section~\ref{sec:task}).

\paragraph{Proposition 1 (stored claims are never evicted).}
For every entity $e$ and cutoff $t<N$, $I_t(e)\subseteq I_{t+1}(e)$. Hence every structured
line rendered for $e$ at $t$ is rendered at $t+1$, possibly with a closing volume added.

\emph{Proof.} $\mathrm{start}(i)\le t$ implies $\mathrm{start}(i)\le t+1$. Structured fields
are rendered deterministically from $I_t(e)$, and an interval that closes after $t$ is shown
as current at $t$ and as a closed interval (``previously'') once its end is reached, so it is
never removed. $\square$

\emph{Consequence.} Under a fixed entity index, a gold fact conveyed by a structured line at
$t$ is conveyed at $t+1$. Losses can come only from prose regenerated per cutoff and from
labelling differences between adjacent cutoffs; \sys loses \res{DROP-anne-B5} conveyed fact
per text (Table~\ref{tab:gains}). The guarantee is conditional on the entity index: with an
index rebuilt per cutoff (Section~\ref{sec:results-scope}) a page's entity can change between
cutoffs.

\paragraph{Proposition 2 (a fixed page budget caps recall).}
Let a system write the page of $e$ in one call under a budget of at most $B$ sentences, each
conveying at most $k$ gold facts. Then at most $kB$ facts are conveyed, so
\begin{equation*}
\mathrm{recall}_t(e)\;\le\;\min\!\Bigl(1,\;\frac{kB}{|\Gt(e)|}\Bigr).
\end{equation*}
\emph{Proof.} Immediate from the count of conveyed facts. $\square$

\emph{Consequence.} For a character whose record keeps growing, the ceiling falls with $t$
once $|\Gt(e)|>kB$; Proposition~1 imposes no such ceiling on rendered claims. The shared
template of B1--B4 allows at most $B=22$ sentences (appearance 2--4, personality 2--6,
history 3--12; Appendix~\ref{app:baselines}), while the eligible checks on the three central
characters grow from 14 at $t=1$ to 38 at $t=5$ (Table~\ref{tab:central-cutoff}). The bound
is consistent with, but does not by itself explain, B1's fall on these characters
(Figure~\ref{fig:accumulation}a): $k$ is not fixed, and the generator also chooses which
facts to keep.

\paragraph{Proposition 3 (a volume filter is blind to modality).}
Let the renderer print $(s,p,o)$ for every claim with $\phi_t(c)=1$, whatever $\mu(c)$.
If some claim has $\mathrm{vol}(c)\le t$, $\mu(c)\neq\textit{asserted}$, and $(s,p,o)$
states a later gold fact $g\in\Ft$, then the page at $t$ discloses $g$.

\emph{Proof.} $\phi_t(c)=1$, so $c$ is rendered, and the rendering states $g$. $\square$

\emph{Consequence.} Filtering by evidence volume is sufficient against disclosure only if
every eligible non-asserted claim is either withheld from deterministic fields or rendered
with its modality and speaker. On the treated text, a neighbour's prediction in book two
was stored as a presumed marriage and printed as a spouse relation at $t=2$ and $t=3$
(\res{M4-annedecon-total-B5}; Section~\ref{sec:results-main}). Prefix-bounded prompting
keeps the passage, with speaker and mood, in the generator's input, and disclosed no
assessed later fact.

\paragraph{Proposition 4 (routing by co-mention loses support).}
Let a scene record $k$ with participants $P_k$ produce a summary line $a$ that is routed to
every pair page $\{u,v\}\subseteq P_k$, and let a line be supported on a pair page only when
$u$ and $v$ share an episode of the scene (the pair rule of Appendix~\ref{app:annotation}).
If the episodes of the scene group the participants into disjoint sets $Q_1,\dots,Q_m$, then
of the $\binom{|P_k|}{2}$ pair pages that receive $a$, at most
$\sum_{j}\binom{|Q_j|}{2}$ can support it.

\emph{Proof.} A pair that shares an episode lies inside one $Q_j$. $\square$

\emph{Consequence.} A line can be true as a statement and unsupported on most of the pages
it reaches, and the unsupported share grows with the participants per record and the
episodes per chapter. In the reported build, scene records list a mean of \res{REV-part}
participants (first build: \res{REV-part-first}); on relationship pages
\res{M3K-anne-B5-rel-atom} of sampled lines are true as statements but
\res{M3K-anne-B5-rel} are supported in page context (Table~\ref{tab:kinds}). Statement
truth and page support should therefore be scored separately
(Appendix~\ref{app:recommendations}).
